% The spacing idiom `\mathopen{}\mathclose{\left( … \right)}` reads as its `\left( … \right)` (OXIDIZED_DESIGN
% #576): `\right` reads its delimiter before it closes its group (tex.web §1191), so the
% `\aftergroup\egroup` of `\mathclose\bgroup\left … \aftergroup\egroup\right` closes the `\mathclose` group after the
% fence; the parser drops the empty OPEN atom and lets the CLOSE wrap holding one fence through
% (`parser.rs::dissolve_class_wraps`). Perl: `not_parsed:>OPEN` for every row (KNOWN_PERL_ERRORS #644).
\documentclass{article}
\usepackage{amsmath}
\begin{document}

% physics2 phy-ab.sty:38-39 `\delopen`/`\delclose` (physics2.sty:83-84); mleftright.sty:246-253, mathtools.sty:831 the
% same (both bound natively here)
\section{The macro idiom}
\makeatletter
\def\delopen{\mathopen{}\mathclose\bgroup\left}
\def\delclose{\aftergroup\egroup\right}
\makeatother
$\delopen( x \delclose)$\par
$f\delopen( x \delclose)$\par
$a\delopen( x+y \delclose) = b$\par
$\delopen| \frac{1}{x} \delclose|$\par
$\delopen[ a, b \delclose)$\par
$\delopen\{ y \middle| y>0 \delclose\}$\par
$\delopen( x \delclose)^2$\par
$\delopen. f(x) \delclose|_{0}^{1}$\par
$\delopen( \delopen( x \delclose) + 1 \delclose)$

% witnesses 2605.22010 (old_macros.tex:119), 2605.21750 (shortcuts-bare.sty:45): the author's `\left`/`\right`
\section{Redefined left and right}
\makeatletter
\begingroup
\let\@left\left \let\@right\right
\def\left{\mathopen{}\mathclose\bgroup\@left}
\def\right{\aftergroup\egroup\@right}
$f\left( x \right)$\par
$\sum_k \left( k + \frac12 \right) x^k$
\endgroup
\makeatother

% A split row's MathFork reads the cells through XMRefs: the dissolved wraps' references follow their children.
\section{A split row}
\begin{equation}
\begin{split}
a &= f\mathopen{}\mathclose{\left( x \right)} \\
  &= g
\end{split}
\end{equation}

% An empty OPEN or CLOSE atom alone is spacing too.
\section{Empty class atoms}
$f\mathopen{}(x)$\par
$a \mathclose{} b$

% Not an idiom: a wrap holding more than one fence keeps its role and reads as written.
\section{A wrap of two fences}
$\mathclose{\left(a\right)+\left(b\right)}$
\end{document}
